\begin{pastexam}[label={px:ust-115-ee-ds-11}]{臺聯大電機類 115 學年度 第 11 題}
（單選；2 分）

To detect whether a cycle exists in a singly linked list, the most famous Floyd's Cycle Detection
Algorithm (the tortoise and the rabbit algorithm) uses two metrics. Assume the distance from the
beginning of the linked list to the start of the cycle is A, and the length of the cycle is L.
Assume the rabbit's speed is twice that of the tortoise. When the tortoise and rabbit meet, assume
the tortoise has traveled a distance B within the cycle. What is the total distance the rabbit has
traveled at this point?

(A)~A + B \qquad (B)~2(A + B) \qquad (C)~A + B + L \qquad (D)~2A + B
\end{pastexam}

\begin{pastsoln}
答案是 (B)。

烏龜和兔子同時從串列開頭出發，走的時間一樣，兔子的速度是烏龜的兩倍，所以兔子走的距離是烏龜的
兩倍（資結書第~5.6~節：做完 $k$ 輪後慢指標在 $x_k$、快指標在 $x_{2k}$）。 烏龜先走 $A$ 到環的
起點，再在環裡走 $B$，共走 $A + B$；兔子走了
\[
2(A + B) .
\]
另一方面，相遇時兔子和烏龜在同一點，兔子比烏龜多繞了整數圈：$2(A + B) = (A + B) + kL$，
$k \ge 1$，也就是 $A + B = kL$。 這個關係正是找環起點的依據（資結書定理~5.4）。

逐項檢查：
\begin{itemize}
\item (A) $A + B$ 是烏龜走的距離，不是兔子的。 錯。
\item (B) 對。
\item (C) $A + B + L$ 是「兔子恰好比烏龜多繞一圈」（$k = 1$）時的值。 $k$ 不一定是 $1$：例如
$A = 5$、$L = 2$，烏龜到達環的起點時兔子已經在環裡繞了好幾圈，相遇時 $A + B = 6 = 3L$，$k = 3$，
兔子走了 $12$，不等於 $A + B + L = 8$。 不是一般的答案。 錯。
\item (D) $2A + B$ 沒有根據；它等於 $2(A+B)$ 只在 $B = 0$ 時成立。 錯。
\end{itemize}
\end{pastsoln}
